%%%%%%%%%%%%%%%%%%%%%%%%%%%%%%%%%%%%%%%%%%%%%%%%%%%%%%%%%%%%
\section{The Modulus Problem for Gradient-Constrained Classes}\label{app:modulus}
%%%%%%%%%%%%%%%%%%%%%%%%%%%%%%%%%%%%%%%%%%%%%%%%%%%%%%%%%%%%

This appendix develops the structure of the modulus of continuity problem for the bivariate gradient-constrained classes $\mathcal{F}(\mathcal{C})$ introduced in \cref{sec:technical:support}, with the goal of characterizing the exact minimax risk constant for the anisotropic rectangle $\mathcal{C}_{\mathrm{rect}} = [-M_x, M_x] \times [-M_z, M_z]$.


\subsection{Setup and the Role of $\delta$}\label{app:modulus:setup}

Consider a fixed design $\{(x_i, z_i)\}_{i=1}^n$ with outcomes $Y_i = f(x_i, z_i) + \varepsilon_i$, $\varepsilon_i \sim \mathcal{N}(0, \sigma^2)$, and target functional $Lf = f(x_0, z_0)$ (evaluation at the cutoff).
The observation operator $K$ maps $f$ to the vector $(f(x_1, z_1), \ldots, f(x_n, z_n))$.

The single-class modulus of continuity for a centrosymmetric class $\mathcal{F}(\mathcal{C})$ is
\begin{equation}\label{eq:app_modulus}
	\omega(\delta;\, \mathcal{F}(\mathcal{C})) = \sup\bigl\{ 2f(x_0, z_0) : \norm{Kf} \leq \delta/2,\; f \in \mathcal{F}(\mathcal{C}) \bigr\}.
\end{equation}

The parameter $\delta$ controls the \emph{noise budget}: the constraint $\norm{Kf} = \sqrt{\sum_i f(x_i, z_i)^2} \leq \delta/2$ requires the function to be small at the observation points, while the objective maximizes $f$ at the target.
A function that achieves a large value of $f(x_0, z_0)$ while remaining small at the observations is one that the data cannot distinguish from the zero function --- the hardest function to estimate.

The modulus $\omega(\delta)$ is increasing and concave in $\delta$: a larger noise budget permits the function to be more visible at the observations, allowing a taller ``tent'' around the target.

\paragraph{Determining $\delta^*$: the bias--variance tradeoff.}
The modulus governs the minimax risk through the following chain.
Fix the function class $\mathcal{F}(\mathcal{C})$, the observation operator $K$ (determined by the design), and the noise level $\sigma$.
For any $\delta > 0$, the worst-case bias of the modulus-optimal affine estimator is $\omega(\delta)/2$ (increasing in $\delta$), while the worst-case standard deviation is $\sigma \delta / \omega(\delta)$ (decreasing in $\delta$ for concave $\omega$).
The optimal tradeoff selects $\delta^*$ as the solution to
\begin{equation}\label{eq:app_delta_star}
	\delta^* = \argmin_{\delta > 0}\; \frac{\omega(\delta)^2}{4} + \frac{\sigma^2 \delta^2}{\omega(\delta)^2},
\end{equation}
or equivalently, the value at which the marginal bias reduction from decreasing $\delta$ equals the marginal variance increase.
By Theorem~3.1 of \textcite{armstrong_optimal_2018}, the minimax risk over $\mathcal{F}(\mathcal{C})$ is governed by $\omega(\delta^*)$.
The practical implication is that once $\omega(\delta;\, \mathcal{F}(\mathcal{C}))$ is known as a function of $\delta$, the noise level $\sigma$ pins down $\delta^*$ and hence the minimax risk and the optimal estimator.

\paragraph{The observation operator $K$ and the sample size $n$.}
The operator $K$ maps a function $f$ to its values at the $n$ design points: $Kf = (f(x_1, z_1), \ldots, f(x_n, z_n)) \in \mathbb{R}^n$.
Thus $K$ is not a primitive of the problem independent of $n$ --- it is \emph{defined} by the design, and adding or removing observations changes $K$.
The modulus $\omega(\delta;\, \mathcal{F}(\mathcal{C}))$ depends on $n$ entirely through $K$: adding an observation point $(x_{n+1}, z_{n+1})$ tightens the data constraint $\norm{Kf} \leq \delta/2$ (the function must now be small at one more point), which shrinks $\omega(\delta)$ for every $\delta$ and hence reduces the minimax risk.
In particular, different designs with the same $n$ yield different moduli --- what matters is not the number of observations per se but their \emph{geometric configuration} relative to the target $(x_0, z_0)$.


\subsection{The Least Favorable Function}\label{app:modulus:lff}

The function achieving the supremum in \eqref{eq:app_modulus} --- the \emph{least favorable function} --- is determined by both the constraint set $\mathcal{C}$ and the design geometry $K$.

\paragraph{Unconstrained case: the Lipschitz tent.}
Ignoring the data constraint $\norm{Kf} \leq \delta/2$, the function maximizing $f(x_0, z_0)$ subject to $\nabla f \in \mathcal{C}$ and $f(x_0, z_0) = 0$ is the support function itself:
\begin{equation}\label{eq:app_tent}
	f^*(x, z) = h_{\mathcal{C}}(x - x_0,\, z - z_0).
\end{equation}
For the rectangle, this is
\begin{equation*}
	f^*_{\mathrm{rect}}(x, z) = M_x\abs{x - x_0} + M_z\abs{z - z_0},
\end{equation*}
which increases linearly from the target in both coordinate directions.
For the disk, $f^*_{\mathrm{disk}}(x, z) = R\sqrt{(x - x_0)^2 + (z - z_0)^2}$, a cone.

\paragraph{Constrained case: the modulated tent.}
The data constraint $\norm{Kf} \leq \delta/2$ forces the function to be small at the observation points.
The unconstrained tent $h_{\mathcal{C}}(x - x_0, z - z_0)$ is large at distant observations, so the actual least favorable function must ``bend down'' --- it rises from the observations toward the target, peaks, and returns.
In the univariate case ($d = 1$), \textcite{armstrong_optimal_2018} show that the least favorable functions are piecewise linear with bandwidths $h_+, h_-$ (their equation~(3)):
\begin{equation*}
	g^*(x) = 1(x \geq 0)(L_0 + b) + C h_+ \cdot k_+(x/h_+) - C h_- \cdot k_-(x/h_-),
\end{equation*}
where $k(u) = \max\{0, 1 - \abs{u}\}$ is the triangular kernel.
The bandwidths are determined by the equal-weighting condition (their equation~(5)), which ensures that positive and negative observations contribute equally to distinguishing the null from the alternative.

In the bivariate case with the rectangle constraint, the least favorable function inherits the separable structure of $h_{\mathrm{rect}}$, but the data constraint couples the coordinates through the design geometry.


\subsection{Derivative Switching Under the Rectangle}\label{app:modulus:switching}

A key structural feature of the rectangle constraint is that the least favorable function can \emph{switch} the sign of $\partial f / \partial z$ across different regions of the design, independently of $\partial f / \partial x$.

Consider the support function tent centered at the target $(x_0, z_0) = (0, 0)$:
\begin{equation*}
	f^*_{\mathrm{rect}}(x, z) = M_x\abs{x} + M_z\abs{z}.
\end{equation*}
This function has gradient
\begin{equation*}
	\nabla f^*_{\mathrm{rect}}(x, z) = \bigl(M_x \sign(x),\; M_z \sign(z)\bigr)
\end{equation*}
at every point where the gradient exists.
The $z$-derivative switches from $+M_z$ above $z = 0$ to $-M_z$ below $z = 0$, while the $x$-derivative independently takes value $+M_x$ or $-M_x$ depending on the sign of $x$.
This switching is feasible because the rectangle is a product set: the pair $(M_x, +M_z)$ and the pair $(M_x, -M_z)$ are both in $[-M_x, M_x] \times [-M_z, M_z]$.

Under the disk constraint, the same switching is feasible only if the switched gradients both lie in $B_2(R)$:
both $(M_x, +M_z)$ and $(M_x, -M_z)$ require $M_x^2 + M_z^2 \leq R^2$.
This is satisfied when $R \geq \sqrt{M_x^2 + M_z^2}$, but at the cost of an inflated bound: the disk must be large enough to contain the corners of the rectangle.

\paragraph{Design-dependent switching pattern.}
The actual least favorable function's switching pattern depends on the design.
For a linear estimator $\hat{L} = \sum_i w_i Y_i$ targeting $f(x_0)$, the worst-case function must simultaneously:
\begin{enumerate}
	\item Maximize $f$ at positive-weight observations (bias contribution $w_i [f(x_i, z_i) - f(x_0)]$ is large positive).
	\item Minimize $f$ at negative-weight observations (bias contribution is large positive when $w_i < 0$ and $f(x_i, z_i) < f(x_0)$).
	\item Satisfy the inter-observation Lipschitz constraints: $\abs{f(x_i, z_i) - f(x_j, z_j)} \leq h_{\mathcal{C}}(x_i - x_j, z_i - z_j)$.
\end{enumerate}

\commentT{The inter-observation Lipschitz constraints (item 3) are the key complication.  Without them, the worst-case bias is simply $\sum_i \abs{w_i} h_{\mathcal{C}}(x_i - x_0, z_i - z_0)$ --- the formula from \cref{prop:wcb}.  But a single function cannot simultaneously achieve $+h_{\mathcal{C}}$ at some observations and $-h_{\mathcal{C}}$ at others if they are too close to each other.  The worst-case function must navigate these pairwise constraints, and the actual worst-case bias is the solution to a constrained optimization (linear program) over the feasible function values $\Delta_i = f(x_i) - f(x_0)$ subject to $\abs{\Delta_i - \Delta_j} \leq h_{\mathcal{C}}(x_i - x_j, z_i - z_j)$ for all pairs.  This is already noted in the proof of \cref{prop:wcb}.}

Whether the derivative switches between $+M_z$ and $-M_z$ at a given observation depends on the relative position of nearby observations in the $z$-direction:
\begin{itemize}
	\item If two cross-stratum observations $(x_i, z_i)$ and $(x_j, z_j)$ with $z_i > z_0 > z_j$ both carry positive weight, the adversary wants $f$ large at both.
	The fastest increase from the target to $z_i > 0$ uses $\partial f / \partial z = +M_z$; the fastest increase to $z_j < 0$ uses $\partial f / \partial z = -M_z$.
	The V-shape $M_z\abs{z}$ achieves both simultaneously.
	\item If one observation carries positive weight and a nearby observation carries negative weight, the adversary wants $f$ large at one and small at the other.
	The derivative sign is chosen to maximize the difference $f(x_i, z_i) - f(x_j, z_j)$, subject to the inter-observation Lipschitz constraint.
\end{itemize}

Under the rectangle, these switches are ``free'': the $x$-derivative can independently take its worst-case value regardless of what $\partial f / \partial z$ is doing.
Under the disk, switching the $z$-derivative from $+M_z$ to $-M_z$ while maintaining a large $x$-derivative requires the gradient to traverse the disk boundary, constraining the achievable $x$-derivative during the transition.


\subsection{The Modulus as a Constrained Program}\label{app:modulus:program}

For a fixed design with $n$ observation points, the modulus problem \eqref{eq:app_modulus} can be written as a finite-dimensional optimization over the function values $\Delta_i = f(x_i, z_i) - f(x_0, z_0)$:
\begin{equation}\label{eq:app_modulus_program}
\begin{aligned}
	\omega(\delta;\, \mathcal{F}(\mathcal{C})) = \sup_{\Delta \in \mathbb{R}^n} \quad & {-2} \cdot 0 \\
	\text{subject to} \quad & \sqrt{\textstyle\sum_i (\Delta_i + f_0)^2} \leq \delta/2, \\
	& \abs{\Delta_i - \Delta_j} \leq h_{\mathcal{C}}(x_i - x_j,\, z_i - z_j), \quad \forall\, i, j, \\
	& \abs{\Delta_i} \leq h_{\mathcal{C}}(x_i - x_0,\, z_i - z_0), \quad \forall\, i,
\end{aligned}
\end{equation}
where the last two families of constraints ensure that the function values are achievable by some $f \in \mathcal{F}(\mathcal{C})$.

\commentT{This formulation needs to be made precise.  The modulus maximizes $2f(x_0)$, which with the normalization $f(x_0) = 0$ becomes: maximize $0$ subject to $f$ being small at the observations --- this is not right.  The correct formulation follows AK's testing interpretation: the modulus is $\sup\{Lg^* - Lf^* : \norm{K(g^* - f^*)} \leq \delta,\; f^*, g^* \in \mathcal{F}\}$, where $f^*$ and $g^*$ are the least favorable pair.  Alternatively, for the single-class modulus with centrosymmetric $\mathcal{F}$: $\omega(\delta) = \sup\{2Lf : \norm{Kf} \leq \delta/2,\; f \in \mathcal{F}\}$.  The finite-dimensional reduction requires showing that the modulus depends on $f$ only through its values at $x_0, x_1, \ldots, x_n$, which holds because the constraints are path-based (Lipschitz along any path connecting observations to the target).  Need to write this out carefully, potentially following Donoho (1994) Section~3 or AK Appendix~A.}

For the rectangle, $h_{\mathrm{rect}}(d_x, d_z) = M_x\abs{d_x} + M_z\abs{d_z}$, so the inter-observation constraints decompose additively:
\begin{equation*}
	\abs{\Delta_i - \Delta_j} \leq M_x\abs{x_i - x_j} + M_z\abs{z_i - z_j}.
\end{equation*}
For the disk, $h_{\mathrm{disk}}(d_x, d_z) = R\sqrt{d_x^2 + d_z^2}$, which does not decompose.

The additive structure of $h_{\mathrm{rect}}$ suggests that the modulus for the rectangle may be bounded or computed using the 1D moduli for $M_x$ and $M_z$ separately, but the design geometry (through the pairwise constraints) couples them.

\commentT{Open problem: can the bivariate modulus for the rectangle be decomposed into (or bounded by) a combination of univariate moduli?  The additive support function $h_{\mathrm{rect}}(v) = M_x\abs{v_1} + M_z\abs{v_2}$ suggests a connection to the product structure, but the data constraint $\norm{Kf} \leq \delta/2$ operates on the joint function values, not on marginals.  If the design is a product grid ($x$-points $\times$ $z$-points), the modulus might factorize.  For general designs, an upper bound via $\omega_{\mathrm{rect}} \leq \omega_{M_x}^{(x)} + \omega_{M_z}^{(z)}$ (sum of marginal moduli) would already give a tighter bound than the disk modulus.  This is the key technical step needed for the exact constant.}


\subsection{Comparison: Rectangle vs.\ Disk Modulus}\label{app:modulus:comparison}

The inclusion $\mathcal{C}_{\mathrm{rect}} \subsetneq \mathcal{C}_{\mathrm{disk}}$ gives $\omega(\delta; \mathcal{F}(\mathcal{C}_{\mathrm{rect}})) \leq \omega(\delta; \mathcal{F}(\mathcal{C}_{\mathrm{disk}}))$ by monotonicity (\cref{thm:monotonicity}).

The gap arises because the disk's least favorable function can concentrate its gradient budget:
\begin{itemize}
	\item Under the disk, the adversary can set $\nabla f = (R, 0)$ --- concentrating the entire budget into $x$.
	Under the rectangle, $\abs{\partial f / \partial x} \leq M_x < R$: the $z$-budget cannot be transferred.
	\item Under the disk, the adversary pays $R\sqrt{\Delta x^2 + \Delta z^2}$ for a diagonal displacement.
	Under the rectangle, the cost is $M_x\abs{\Delta x} + M_z\abs{\Delta z}$ --- the coordinates contribute additively, and the $z$-cost is ``stranded'' for horizontal displacements.
\end{itemize}

The exact ratio of moduli, and hence of minimax risks, depends on the design through the displacement vectors $\{(x_i - x_0, z_i - z_0)\}_{i=1}^n$.
By \cref{lem:ratio_range}, the pointwise support function ratio satisfies $1 \leq r(v) \leq R / \min(M_x, M_z)$, which bounds the modulus ratio.
Computing the exact ratio for a specific design is an open problem (\cref{sec:technical:open}).


\subsection{The Separability Argument for Second-Order Classes}\label{app:modulus:hessian_norms}

The separability trilemma of \cref{prop:coupling_agnostic} applies not only to first-order (gradient) constraints but extends to second-order (Hessian) constraints, clarifying the norm choice in \textcite{imbens_optimized_2019}.

\paragraph{Entry-wise bounds as the separable class.}
For a bivariate function $f(x,z)$, the Hessian is the symmetric matrix
\begin{equation*}
	\nabla^2 f = \begin{pmatrix} \partial^2 f / \partial x^2 & \partial^2 f / \partial x \partial z \\ \partial^2 f / \partial x \partial z & \partial^2 f / \partial z^2 \end{pmatrix},
\end{equation*}
with three free parameters: the diagonal entries $a = \partial^2 f / \partial x^2$, $c = \partial^2 f / \partial z^2$ and the cross-partial $b = \partial^2 f / \partial x \partial z$.
If a researcher has separate bounds $\abs{a} \leq B_1$, $\abs{b} \leq B_2$, $\abs{c} \leq B_3$ on each second derivative, the constraint set in $(a,b,c)$-space is the cuboid $[-B_1, B_1] \times [-B_2, B_2] \times [-B_3, B_3]$.
This cuboid is the \emph{separable} constraint for second derivatives: it bounds each entry independently without imposing any relationship between them --- the direct analog of the rectangle for gradients.

\paragraph{The operator norm imposes coupling.}
\textcite{imbens_optimized_2019} bound the operator norm $\norm{\nabla^2 f}_{\mathrm{op}} = \max(\abs{\lambda_1}, \abs{\lambda_2}) \leq B$, where the eigenvalues are
\begin{equation*}
	\lambda_{\pm} = \frac{a + c}{2} \pm \sqrt{\left(\frac{a-c}{2}\right)^2 + b^2}.
\end{equation*}
The eigenvalues are nonlinear functions of the entries, coupling $a$, $b$, and $c$ jointly.
The key inequalities between the operator norm and the entry-wise maximum are
\begin{equation}\label{eq:entry_vs_op}
	\max_{i,j} \abs{H_{ij}} \;\leq\; \norm{H}_{\mathrm{op}} \;\leq\; n \cdot \max_{i,j} \abs{H_{ij}},
\end{equation}
where $n = 2$ in the bivariate case.
The left inequality shows that the operator-norm ball $\{\norm{H}_{\mathrm{op}} \leq B\}$ is contained in the entry-wise cube $\{\max \abs{H_{ij}} \leq B\}$: the operator norm is the stricter constraint.
The containment is strict: the matrix $\bigl(\begin{smallmatrix} B & B \\ B & B \end{smallmatrix}\bigr)$ has all entries equal to $B$ but $\lambda_{\max} = 2B$, so it lies in the cube but outside the operator-norm ball.

The operator norm constrains the entries \emph{jointly} through the eigenvalue structure: a large diagonal entry forces the off-diagonal to be small, and vice versa.
This is coupling in the sense of \cref{def:coupling} --- the second-order analog of the disk constraining $(\partial f / \partial x,\, \partial f / \partial z)$ jointly.

\paragraph{The trilemma extends.}
By the logic of \cref{def:coupling_agnostic}, coupling-agnosticism requires the constraint set to contain the separable class (the cuboid).
For the operator-norm ball, containing the cube $[-B, B]^3$ requires radius at least $nB$ (to accommodate the corner $(B, B, B)$, which has $\norm{H}_{\mathrm{op}} = nB$).
But an operator-norm ball of radius $nB$ admits entry-wise values up to $nB$ by the left inequality in \eqref{eq:entry_vs_op} --- inflating the primitive bounds by a factor of $n$.

The same trilemma from the first-order case applies to second-order classes:
\begin{itemize}
	\item The operator-norm ball at radius $B$ is a strict subset of the entry-wise cube at radius $B$: it imposes coupling on the second derivatives.
	\item To be coupling-agnostic (contain the cube), the operator-norm ball must inflate to radius $nB$, exceeding the primitive entry-wise bounds.
	\item The entry-wise cuboid is the unique separable class that is simultaneously coupling-agnostic and respects the primitive bounds on each second derivative.
\end{itemize}

\paragraph{Eigenvalue space clarifies the geometry.}
In eigenvalue coordinates $(\lambda_1, \lambda_2)$, the standard matrix norms take the same forms as vector norms:
\begin{center}
\begin{tabular}{lll}
	Norm & Eigenvalue constraint & Shape in $(\lambda_1, \lambda_2)$ \\
	\hline
	Operator ($\norm{\cdot}_{\mathrm{op}}$) & $\max(\abs{\lambda_1}, \abs{\lambda_2}) \leq L$ & Square ($\ell^\infty$-ball) \\
	Frobenius ($\norm{\cdot}_F$) & $\sqrt{\lambda_1^2 + \lambda_2^2} \leq L$ & Disk ($\ell^2$-ball) \\
	Nuclear ($\norm{\cdot}_*$) & $\abs{\lambda_1} + \abs{\lambda_2} \leq L$ & Diamond ($\ell^1$-ball)
\end{tabular}
\end{center}
The Frobenius disk is inscribed in the operator-norm square, just as the $\ell^2$-disk is inscribed in the $\ell^\infty$-square for gradients.
The entry-wise cube, mapped to eigenvalue space, produces a \emph{bulged} region that strictly contains the operator-norm square at the same radius: the off-diagonal entry $b$ contributes to eigenvalue spread through $\sqrt{((a-c)/2)^2 + b^2}$ without being directly constrained by the eigenvalue bound.

\commentT{The parallel between first-order and second-order separability is structurally clean and suggests a general principle.  For gradients: the rectangle (product of marginal bounds) is separable, the disk (joint norm) imposes coupling.  For Hessians: the entry-wise cuboid is separable, the operator-norm ball imposes coupling.  In both cases, the coupling-agnostic class is the entry-wise constraint, and any joint norm constraint must either impose coupling or inflate the primitive bounds.  The inflation factor scales with dimension $n$.  This extends to any order of derivatives: for a $k$-th order derivative tensor, the entry-wise bound on its unique elements is the separable class, and any rotationally-invariant norm constraint imposes coupling between them.}
